\documentclass[tikz]{standalone}

\usepackage{tikz}
\usetikzlibrary{arrows.meta,positioning,fit,shapes.misc}

\begin{document}

\begin{tikzpicture}[scale=0.85, every node/.style={font=\small}]
  %----- Styles -----
  \tikzset{
    axis/.style={->, thick},
    gwc/.style={dashed, very thick},
    gridline/.style={gray!40},
    rhfill/.style={fill=black!15, opacity=0.5},
    rhborder/.style={black!50, thick},
    lwcfill/.style={fill=black!30, opacity=0.55},
    point/.style={circle, fill=black, inner sep=1.2pt},
    legendbox/.style={draw=black!50, rounded corners=2pt, fill=white, line width=0.3pt},
    rect/.style={black!100, very thick},
  }

  %----- Axes -----
  \draw[axis] (0,0) -- (9,0) node[below right] {$E_1$};
  \draw[axis] (0,0) -- (0,6.8) node[above left]  {$E_2$};

  %----- Global GWC rectangle (W_1^gwc, W_2^gwc) -----
  \def\Wone{6.6}
  \def\Wtwo{4.7}
  \draw[gwc] (0,0) rectangle (\Wone,\Wtwo);

  %----- Partition grid (order statistics) -----
  % vertical grid lines (order stats of E1)
  \foreach \x/\lab in {1.6/$E^{(1)}_1$, 3/$E^{(2)}_1$, 4/$E^{(3)}_1$} {
    \draw[gridline] (\x,0) -- (\x,6.3);
    \node[gray!60] at (\x,-0.35) {\lab};
  }
  % horizontal grid lines (order stats of E2)
  \foreach \y/\lab in {0.9/$E^{(1)}_2$, 1.8/$E^{(2)}_2$, 2.7/$E^{(3)}_2$} {
    \draw[gridline] (0,\y) -- (8.2,\y);
    \node[gray!60,rotate=90] at (-0.35,\y) {\lab};
  }
  \node[gray!60] at (0,-0.35) {$0$};

  %----- (Optional) visually frame one R_h cell (border only) -----
  \def\xL{4} \def\xR{5.5} \def\yB{1.8} \def\yT{2.7}
  %----- A few example LWC subregions E^{lwc,h} (asymmetric slivers inside their R_h) -----
  % Example 1 inside R_(3,3)
  \fill[lwcfill] (\xL,\yB) -- (\xR-0.25,\yB) -- (\xR-0.25,\yT) -- (\xL,\yT) -- cycle;
  \fill[lwcfill] (\xL,\yB-0.9) -- (\xR-0.5,\yB-0.9) -- (\xR-0.5,\yT-0.9) -- (\xL,\yT-0.9) -- cycle;
  \fill[lwcfill] (\xL,0) -- (\xR-0.5,0) -- (\xR-0.5,\yT-1.8) -- (\xL,\yT-1.8) -- cycle;

  % Example 2 inside R_(5,2): [4.8,6.0] x [0.9,1.8]
  \def\xLtwo{0} \def\xRtwo{4} \def\yBtwo{0} \def\yTtwo{2.7}
  \fill[lwcfill] (\xLtwo,\yBtwo) -- (\xRtwo,\yBtwo) -- (\xRtwo,\yTtwo) -- (\xLtwo,\yTtwo) -- cycle;

  % Example 3 inside R_(2,5): [1.2,2.4] x [3.6,4.2]
  \def\xLthr{0} \def\xRthr{1.6} \def\yBthr{2.7} \def\yTthr{3.2}
  \fill[lwcfill] (\xLthr,\yBthr) rectangle (\xRthr,\yTthr+0.2);
  \fill[lwcfill] (\xLthr+1.6,\yBthr) rectangle (\xRthr+1.2,\yTthr+0.2);
  \fill[lwcfill] (\xLthr+3,\yBthr) rectangle (\xRthr+2.4,\yTthr+0.4);
  \fill[lwcfill] (\xLthr+4,\yBthr) rectangle (\xRthr+3.4,\yTthr+0.4);

\def\vone{5.25}
  \def\vtwo{3.6}
  \draw[rect] (0,0) rectangle (\vone,\vtwo);  


  %----- Candidate test residual point E^{n+1} -----
  %\fill[black] ($( \xL+0.55, \yB+0.55 )$) circle (1.2pt) node[above right] {$E^{n+1}$};

  %----- (Optional) Enclosing rectangle for the shortcut (dash-dot); uncomment in shortcut section -----
  %\draw[dash dot, ultra thick] (0,0) rectangle (5.6,4.0);
  %\node[anchor=south west] at (0.15,4.0) {enclosing shortcut rectangle};

  %----- Compact legend -----
  \begin{scope}[shift={(10.0,4.0)}]
    \node[legendbox, align=left, inner sep=4pt] (L) {
      \begin{tikzpicture}[baseline=(current bounding box.center)]
        \draw[gwc] (0,0) -- (0.8,0); \node[right] at (1.05,0) {$\hat{C}^{\mathrm{gwc}}(X^{n+1})$};
        \fill[lwcfill] (0,-0.9) rectangle +(0.8,0.18); \node[right] at (1.05,-0.81) {$\hat{C}^{\mathrm{lwc}}(X^{n+1}, \mathcal{R})$};
        \draw[rect] (0,-1.5) -- (0.8,-1.5); \node[right] at (1.05,-1.5) {$\hat{C}^\mathrm{lwc}(X^{n+1})$};
      \end{tikzpicture}
    };
  \end{scope}

\end{tikzpicture}
\end{document}


%%% Local Variables:
%%% mode: latex
%%% TeX-master: t
%%% End:
